\label{ch:inference-process}

This chapter documents the statistical inference machinery: the
normalization-calibration step every weighted-MC fit here needs before it
can start (\cref{sec:calibration}); the Poisson fit statistic and
gradient-based procedure, shared unchanged with the rest of
\texttt{tpeanuts} (\cref{sec:fit-procedure-icecube}); this document's own
dedicated test suite, added specifically to cover this project's only
weighted-MC detector (\cref{sec:test-suite}); and the muon-background
normalization audit this document's own working session carried out
(\cref{sec:munorm-audit}).

\section{Normalization calibration}
\label{sec:calibration}

\begin{theorybox}
Since the release's \code{weight} column has no independently usable
absolute livetime (\cref{sec:free-normalization-icecube}), a real fit
cannot simply start at $\nu_{\rm norm}=1$: the neutrino signal there is
$\sim$0.2\% of the total rate, so the Poisson gradient for
$\theta_{23}$/$\Delta m^2_{3l}$ is numerically negligible and a
gradient-based optimizer's own step-size heuristics can take a
catastrophic first step. A physically sensible starting guess is computed
once, at NuFIT~6.1 defaults \emph{before any fit}, from the real observed
total and real MC muon-background total:
\begin{equation}
  \keyresult{\nu_{\rm norm}^{(0)} = \frac{n^{\rm obs}_{\rm IceCube} - N^{\mu,{\rm MC}}_{\rm IceCube}}
    {N^\nu_{\tt tpeanuts}(\theta_{\rm NuFIT})}, \qquad \mu_{\rm norm}^{(0)}=1,}
  \label{eq:calibration}
\end{equation}
with $N^\nu_{\tt tpeanuts}(\theta_{\rm NuFIT})$ the \texttt{tpeanuts}-simulated
neutrino total at $\nu_{\rm norm}=1$, oscillation parameters fixed at
NuFIT. This is a calibration of scale only, not a measurement or a fit --
the oscillation parameters are not touched, and \cref{sec:fit-procedure-icecube}'s
real fit re-optimizes $\nu_{\rm norm}$/$\mu_{\rm norm}$ freely rather than
leaving them at this starting point.
\end{theorybox}

\begin{resultbox}
At NuFIT~6.1 defaults, $N^\nu_{\tt tpeanuts}(\theta_{\rm NuFIT})=1.0487$ at
$\nu_{\rm norm}=1$ (real, downsampled MC), giving
$\nu_{\rm norm}^{(0)}=20{,}408.5$ -- by construction, the calibrated
starting point's total prediction ($21{,}914.0$) exactly matches the real
observed total.
\end{resultbox}

\section{Poisson fit statistic and gradient-based procedure}
\label{sec:fit-procedure-icecube}

\modulehead{inference/likelihood.py, inference/fit.py}{\pyfunc{poisson\_nll}, \pyfunc{fit\_lbfgs}, shared unchanged with every other detector in this project.}

\begin{theorybox}
Real observed counts $n_k^{\rm obs}$ are compared to
\labelcref{eq:full-prediction} via the Poisson (Baker-Cousins) likelihood
ratio \parencite{BakerCousins1984},
\begin{equation}
  \keyresult{-2\ln L = 2\sum_k\Bigl[N_k - n_k^{\rm obs} + n_k^{\rm obs}\ln\frac{n_k^{\rm obs}}{N_k}\Bigr],}
  \label{eq:poisson-nll-icecube}
\end{equation}
minimized over the full 4-parameter $\theta_{\rm fit}$ with
\texttt{torch.optim.LBFGS}, using the exact autograd gradient through
every step of \cref{ch:detector-implementation} (no finite differences).
The marginal uncertainty on each parameter is the Laplace approximation,
$\Sigma=\mathcal I^{-1}$ with $\mathcal I=\tfrac12\nabla^2_\theta(-2\ln L)$
the autograd Hessian at the optimum -- identical machinery to Daya Bay's
and SNO-II's own fits in this project.
\end{theorybox}

\begin{implbox}
\pyfunc{fit\_lbfgs(model, theta0, value, likelihood="poisson",
max\_iter=...)} returns a \pyclass{FitResult}. Each LBFGS closure
evaluation re-runs the full real-MC reweighting ($\sim$40{,}000 downsampled
events across 4 channels) -- $\sim$4~s/evaluation, so a real fit takes a
few minutes wall time (\cref{ch:results}).
\end{implbox}

\section{This document's own test suite}
\label{sec:test-suite}

\begin{theorybox}
Before this document's working session, \texttt{detector.icecube} had no
dedicated tests: every other detector package in this project already had
its own gradient/conservation/closure test coverage, but IceCube's
event-by-event weighted-MC reweighting -- the only fold of its kind in
this project -- did not. Five tests were added, each targeting a specific
correctness property of \cref{ch:detector-implementation}'s own formulas:
\begin{enumerate}
  \item \textbf{Gradient flow}: autograd reaches all four free parameters
    with finite, non-zero gradients.
  \item \textbf{Reweighting conservation}: the per-event $\to$ per-bin
    histogram (\labelcref{eq:bin-sum}) neither loses nor duplicates
    events, verified against an independently computed direct per-event
    sum at a trivial (all-ones) hypersurface correction.
  \item \textbf{Full-flavour-sum conservation}: summing the reweighted
    prediction over all 3 possible final-state flavours recovers the
    unoscillated rate $w_i(\Phi_e+\Phi_\mu)$, a direct consequence of the
    real Earth-matter transition matrix's own unitarity.
  \item \textbf{Hypersurface interpolation}: exact at tabulated
    $\Delta m^2_{31}$ slices, linear at an interior point, clamped outside
    the tabulated range.
  \item \textbf{Synthetic closure}: fits synthetic Poisson counts
    generated at a known injected $(\theta_{23},\Delta m^2_{3l},\nu_{\rm norm},
    \mu_{\rm norm})$ and confirms recovery.
\end{enumerate}
\end{theorybox}

\begin{notebox}
\textbf{Two genuine numerical traps, found and documented while writing
test 5, not anticipated going in.} An injected/starting $\nu_{\rm norm}$
near 1 (rather than the real $\sim2\times10^4$ scale,
\cref{sec:calibration}) makes the neutrino signal numerically invisible
against the muon background, collapsing the $\theta_{23}$/$\Delta
m^2_{3l}$ gradient; combined with \pyfunc{minimize\_lbfgs}'s own
per-parameter $1/|\nabla|$ rescaling heuristic, this produced a
catastrophic first optimizer step (observed: $\theta_{23}$ jumping to
$-207$~radians on the very first closure evaluation) before the test was
corrected to inject/start at the real, calibrated $\nu_{\rm norm}$ scale.
Separately, a $\theta_{23}$ offset crossing $45^\circ$ (the maximal-mixing
octant boundary, close to the real NuFIT starting value of $43.29^\circ$)
can converge to the physically real, but confusingly different-looking,
octant-degenerate solution -- worked around by keeping the injected
offset small and same-octant, not by loosening the test's own tolerance.
\end{notebox}

\begin{resultbox}
All 5 tests pass; the full project regression suite (848 tests, at the
point this test file was added) shows no regressions.
\end{resultbox}

\section{The muon-background normalization audit}
\label{sec:munorm-audit}

\begin{theorybox}
\Cref{ch:results} reports a real, quantitative discrepancy between this
project's own fitted $\mu_{\rm norm}$ and IceCube's own published
muon-background scale (Table IV). This document's own working session
investigated two candidate explanations, in the same primary-source-verified
spirit as this project's own established practice:
\begin{enumerate}
  \item \textbf{Earth-matter unitarity leak.}
    \pyfunc{earth\_probability\_transition} defaults to
    \code{reunitarize=False} (\texttt{config.default.earth\_reunitarize}),
    so the raw perturbative Earth evolution operator is only unitary up to
    its own real numerical precision -- empirically up to $\sim$0.7\%
    per-event row-sum deviation at IceCube's real energy/baseline range.
    \textbf{Ruled out}, quantitatively: forcing \code{reunitarize=True} at
    the published best-fit point shifts the total predicted rate by only
    $\sim$0.002\% (21{,}913.8 $\to$ 21{,}913.4 of 21{,}914) -- three orders
    of magnitude too small to explain the real gap in \cref{ch:results}.
  \item \textbf{Fixed vs.\ free detector systematics.} The 5 real
    detector-systematics parameters are held fixed at IceCube's own
    published \emph{joint}-fit best-fit point (\cref{sec:flux-hypersurface}),
    which is itself correlated with IceCube's own published $\mu_{\rm norm}$
    in a $\sim$40-parameter joint fit with Gaussian priors on every
    nuisance. This project's own fit has only 4 free parameters, no
    priors, and no floating systematics -- so its own $\mu_{\rm norm}$ has
    to absorb whatever correlated normalization the official fit
    distributed across the (here fixed) systematics. \textbf{Identified as
    the more likely explanation}, and quantified directly: the real
    hypersurface correction's own mean value, best-fit vs.\ nominal
    systematics, ranges from $+2\%$ to $+12\%$ across the four channels
    -- not negligible, and the natural size of correlated normalization
    slack this project's reduced, prior-free fit would need to absorb
    through $\mu_{\rm norm}$ alone.
\end{enumerate}
Floating the 5 systematics as free nuisances with Gaussian priors (the
same \pyfunc{gaussian\_prior\_penalty} machinery already used elsewhere in
this project) would test this hypothesis directly, but was not attempted
in this document -- an explicitly deferred next step
(\cref{ch:results}).
\end{theorybox}
